% =============================================================================
\section{Limitations}
\label{sec:limitations}
% =============================================================================

Our approach has several limitations that should be acknowledged.

\textbf{Statistical fragility on small TREC DL test sets}: With 43 (DL19) and 54 (DL20) queries, only one DualEnd configuration (Qwen3-4B DL19) survives Bonferroni correction with $p_\text{Bonf}=0.010$. The remaining 13 positive DualEnd configurations are directionally consistent but cannot be claimed as significant on these query budgets. By contrast, the negative results---6 Bonferroni-significant BottomUp losses and 3 BiDir losses---are statistically more robust. We have therefore framed the headline contribution as ``directional asymmetry'' rather than ``DualEnd is universally significantly better.''

\textbf{Efficiency of bottom-up for top-$k$}: \bottomup{} requires $\sim$$n-1$ extractions for top-$k$ ranking, making it $\sim$$3.5\times$ more expensive than \topdown{}-Heap when $k=10$, $n=100$ (and \bottomup{}-Bubble is more than $20\times$ more expensive). This limits its use as a standalone method; it is primarily valuable as a diagnostic tool for the directional asymmetry argument.

\textbf{DualEnd cost}: \dualend{}-Cocktail uses about $7\times$ more comparisons than \topdown{}-Heap and about $1.7\times$ more than \topdown{}-Bubble. The token-level overhead is similar. We do not claim end-to-end efficiency gains; we claim more useful bits per call. The Pareto analysis in \S\ref{sec:pareto} shows that the global mean frontier contains only \topdown{}-Heap, \topdown{}-Bubble, and \dualend{}-Cocktail, and this is the right way to read the cost story.

\textbf{Likelihood for worst-related signals remains heuristic}: For Flan-T5 we use direct logits over single-character labels. For Qwen3/Qwen3.5 we use a teacher-forced short-answer scheme (\texttt{Passage A}, etc.) that is robust to causal tokenization. For \emph{dual-end} likelihood we reuse the best-only distribution and take its maximum and minimum as the best/worst outputs. This is useful as an efficiency probe and matches the heuristic Flan-T5 path, but it is not an exact likelihood model of the full ``Best: X, Worst: Y'' string. The same caveat applies to the routed joint-signal refinements in \S\ref{sec:refinement} whenever they enter their joint-signal path under \texttt{--scoring likelihood}.

\textbf{Parsing fragility}: Dual-end generation requires the LLM to produce structured output with two labels. While our cascading parser handles most formats, the success rate varies across models. Models that tend to produce verbose explanations rather than concise labels may have lower parsing success rates. The worst case in our experiments is Qwen3-14B on DL 2020 with a 0.9\% partial parse rate (handled by heuristic defaults).

\textbf{Model dependence}: Our experiments cover three architecture families (encoder-decoder, decoder-only, and hybrid causal) and nine models ranging from 780M to 27B parameters. Results may differ for larger models (e.g., 70B+) or API-based models (GPT-4, Claude), which may have stronger format-following abilities and different position bias patterns.

\textbf{Dataset scope}: Our headline results come from TREC DL 2019/2020. We are also running a representative BEIR subset on Flan-T5-XL, Qwen3-8B, and Qwen3.5-9B; those runs are still completing and will be folded into the camera-ready version. All current benchmarks use relatively short passages, and performance on longer documents or domain-specific corpora with very long documents remains to be validated.

\textbf{Refinement variants are implemented but not yet fully evaluated}: The Selective DualEnd, bias-aware DualEnd, and same-call worst-signal regularization variants in \S\ref{sec:refinement} are runnable code paths in our \texttt{llm-rankers} fork, with their own logged counters. We use them in this paper as a motivated next-step package rather than as a primary empirical contribution; their full evaluation is reserved for follow-up work.


% =============================================================================
\section{Conclusion}
\label{sec:conclusion}
% =============================================================================

We introduced three bidirectional strategies for LLM-based setwise document ranking that challenge the convention of extracting only ``best'' information from each LLM comparison. Our experiments with nine LLMs across three architecture families yield three findings.

First, \textbf{the DualEnd family is the strongest overall}. \dualend{}-Cocktail achieves the highest average NDCG@10 across all 9 models on both DL 2019 (.7134) and DL 2020 (.6791), and DualEnd methods (Cocktail or Selection) win in 14 of 18 model--dataset configurations. The per-query significance picture is more conservative: only one of those gains remains Bonferroni-significant. By asking the LLM for both the best and worst document per call, dual-end methods narrow the ranking from both ends, and an unexpected benefit is \emph{implicit debiasing}: dual-end prompting produces more uniform position selection patterns than single-direction methods, partially mitigating the strong recency bias observed in standalone best- and worst-selection.

Second, \textbf{bottom-up ranking is substantially weaker than top-down}, refuting the hypothesis that LLMs are better at identifying irrelevant documents. The gap is particularly severe for smaller models (Flan-T5-Large: .289 vs.\ .654 NDCG@10) and narrows but does not vanish at scale. Bonferroni-significant losses appear in 6 of 18 configurations. Position bias analysis reveals a likely mechanism: models exhibit extreme recency bias when selecting ``worst'' (position D selected 40--63\% of the time), effectively defaulting to the last passage.

Third, \textbf{bidirectional ensemble does not outperform the top-down baseline}, indicating that simply fusing top-down and bottom-up rankings cannot compensate for the weakness of the bottom-up component. Even at $\alpha=0.9$ (90\% TopDown weight), fusion underperforms pure TopDown, and BiDir produces 3 Bonferroni-significant losses. The intuitive assumption that directional fusion would parallel the success of permutation voting~\citep{tang2024found} does not hold because the directions are not equally reliable.

The strongest scientific contribution is therefore one of \emph{directional asymmetry}: worst-selection alone is not the dual of best-selection, and the gap between standalone worst-selection and joint best-worst elicitation is mechanistically meaningful. The strongest practical contribution is the empirical map that says which sorting algorithm and which selection direction works for which model family, and the quality--cost Pareto frontier that pinpoints the missing intermediate point between \topdown{}-Bubble and \dualend{}-Cocktail.

Our findings open several avenues for future work: (1) full empirical evaluation of the Selective, bias-aware, and same-call regularized DualEnd refinement variants, all of which are already implemented and instrumented in our fork; (2) extending dual-end to extract a \emph{full ordering} of the candidate set, effectively bridging setwise and listwise paradigms; (3) investigating whether reasoning-enhanced rerankers~\citep{zhuang2025rankr1} benefit from dual-end prompting; (4) benchmarking the new Qwen/Qwen3.5 likelihood follow-up against the generation results to quantify the cost saved and the quality lost; and (5) understanding why worst-selection is harder for LLMs and whether this can be mitigated through chain-of-thought reasoning.
